\section{Objetivos}

\subsection{Objetivo general}

Desarrollar un sistema distribuido para gestionar los documentos de las
pre-sustentaciones de trabajos de titulación en la Universidad Técnica Estatal
de Quevedo, desde la entrega de la solicitud hasta la generación del acta
de predefensa y la consolidación de la nota final.

\textbf{Indicador.} Porcentaje de casos de prueba de extremo a extremo que
completan el flujo previsto, desde la entrega de documentos hasta la generación
del acta de predefensa y la consolidación de la nota final. \textbf{Meta.} Al menos el 85\%.

\subsection{Objetivos específicos}

\begin{enumerate}
    \item Implementar la entrega del informe de predefensa, el reporte de
    coincidencias de URKUND y el informe corregido cuando haya observaciones.
    La entrega del anteproyecto corresponde a una etapa previa.

    \textbf{Indicador.} Porcentaje de casos de prueba en los que el sistema
    acepta solicitudes con documentos válidos y detecta los documentos
    obligatorios faltantes. \textbf{Meta.} Al menos el 90\%.

    \item Facilitar la revisión de los documentos por parte de coordinación y
    permitir que el estudiante consulte el estado y las observaciones de su
    solicitud.

    \textbf{Indicador.} Porcentaje de solicitudes de prueba en las que el
    estado y las observaciones consultadas por el estudiante coinciden con lo
    registrado por coordinación. \textbf{Meta.} Al menos el 90\%.

    \item Registrar la calificación del tribunal según su rúbrica y
    consolidar la nota final con el 40\% del tribunal y el 60\% del docente
    de Desarrollo de Titulación II.
    En el proceso actual, el docente registra la nota final en el SGA.

    \textbf{Indicador.} Porcentaje de casos de prueba en los que la
    calificación del tribunal queda vinculada a la solicitud correcta y
    la nota final coincide con el cálculo esperado. \textbf{Meta.} Al menos el 95\%.

    \item Generar el acta de predefensa con las observaciones, la calificación
    vigente del tribunal y los espacios de firma del presidente y los dos
    miembros. Bloquear la generación si la calificación vigente no se
    puede verificar.

    \textbf{Indicador.} Porcentaje de casos de prueba que generan el acta
    correcta cuando la nota es verificable o bloquean la generación cuando
    no se puede verificar. \textbf{Meta.} Al menos el 95\%.
\end{enumerate}
